% ============================================================================
% 02_math_chem.tex  —  mathematics, units and chemical notation
%
% Units per docs/00_SPEC.md §12 (Bible §12):
%   capacities in mg/g AND mmol/g   ·   energies in kJ/mol throughout
%   never kcal/mol from software output
% ============================================================================

% yhsa-jinst loads newtxmath, which already defines \Bbbk. amssymb defines it
% again and errors. Release it before loading amssymb rather than dropping
% amssymb, which supplies symbols the report needs.
\usepackage{amsmath}
\let\Bbbk\relax
\usepackage{amssymb}
\usepackage{siunitx}

\sisetup{
  detect-all,
  separate-uncertainty = true,      % 40.1 +- 2.0, not 40.1(20)
  multi-part-units     = single,
  range-units          = single,
  per-mode             = symbol,
  exponent-product     = \times,
  group-digits         = integer,
  group-minimum-digits = 4,
  uncertainty-mode     = separate,
}

% Project units, declared once so they cannot drift between sections.
\DeclareSIUnit{\molar}{\mole\per\litre}
\DeclareSIUnit{\bedvolume}{BV}
\newcommand{\mgg}{\unit{\milli\gram\per\gram}}
\newcommand{\mmolg}{\unit{\milli\mole\per\gram}}
\newcommand{\kJmol}{\unit{\kilo\joule\per\mole}}
\newcommand{\JmolK}{\unit{\joule\per\mole\per\kelvin}}
\newcommand{\mgL}{\unit{\milli\gram\per\litre}}
\newcommand{\mmolL}{\unit{\milli\mole\per\litre}}
\newcommand{\Lg}{\unit{\litre\per\gram}}

% siunitx defines \num. This project reserves \num{} for canonical-number lookup
% (see 06_macros.tex), because NO numeric literal may be hard-coded in any .tex
% file. siunitx's own numeric formatter is preserved as \sinum for internal use.
\let\sinum\num
\let\num\relax

% pH is set upright with a lowercase p, per IUPAC. \pH is \ensuremath-wrapped so
% it works in both modes; a subscript must therefore NOT be attached to it from
% text mode (\pH_{...} puts the subscript outside maths). Use \pHpzc instead.
\newcommand{\pH}{\ensuremath{\mathrm{pH}}}
\newcommand{\pHpzc}{\ensuremath{\mathrm{pH}_{\mathrm{PZC}}}}

\usepackage[version=4]{mhchem}     % \ce{Pb^2+}, \ce{PbSO4}
\mhchemoptions{textfontcommand=\rmfamily}
